\FloatBarrier
\appendix
\section{Instruction encoding and register map}\label{app:isa}

\begin{table}[h]
\centering\small
\caption{AVISA field encoding.}
\begin{tabular}{@{}llll@{}}
\toprule
Field & Bits & Position & Meaning\\
\midrule
opcode & 8 & 63--56 & operation (Table~\ref{tab:isa})\\
precision & 3 & 55--53 & 0 fp32, 1 fp16, 2 bf16, 3 fp8, 4 adaptive\\
vector mask & 3 & 52--50 & lane-group predicate (8 groups of 32 lanes)\\
rd & 5 & 49--45 & destination tile register\\
rs1 & 5 & 44--40 & first source tile register\\
rs2 & 5 & 39--35 & second source tile register\\
descriptor & 24 & 34--11 & stencil coefficients, halo offset, stride\\
member base & 8 & 10--3 & first ensemble member in the lane group\\
flags & 3 & 2--0 & Kahan, saturate, audit\\
\bottomrule
\end{tabular}
\end{table}

The tile register file has 32 registers of 256 lanes $\times$ 32 bits. The descriptor is interpreted
per opcode: for \code{STENCIL3D} it packs three 6-bit coefficients and a 6-bit halo mask; for
\code{DA.GAIN} it selects a localisation radius and a symmetry flag. There is no general addressing
mode; all memory references are offsets into the tile's 256 KB SRAM.

\section{Model parameters}\label{app:params}

\begin{table}[h]
\centering\small
\caption{Parameters of the pre-silicon model. The ``range'' column is the interval sampled in the
Monte-Carlo study.}
\begin{tabular}{@{}lrll@{}}
\toprule
Parameter & Value & Range & Source\\
\midrule
fp32 FMA energy & 0.42 pJ & $\pm$25\% & Horowitz 2014, scaled to 3 nm / 0.75 V\\
fp16 FMA energy & 0.24 pJ & $\pm$25\% & same, half-rate\\
fp8 MAC energy & 0.14 pJ & $\pm$25\% & same\\
SRAM read / write & 0.30 / 0.36 pJ per word & $\pm$25\% & published 3 nm SRAM\\
L3 (stacked) & 0.40 pJ/B & $\pm$25\% & hybrid-bonded SRAM estimate\\
HBM3E & 16.0 pJ/B & 9--24 & HBM3E datasheet, published PHY energy\\
Clock and PLL & 2.5 W & & design budget\\
Leakage & 4.0 W at 0.75 V, 85 $^\circ$C & & library estimate\\
GPU stencil efficiency & 0.30 & 0.15--0.60 & WRF / COSMO GPU studies\\
GPU power at load & 263 W on W1 & $\pm$30\% & H100 data sheet, load model\\
H100 peak fp32 / HBM & 67 TFLOPS / 3.35 TB/s & & data sheet\\
A100 peak fp32 / HBM & 19.5 TFLOPS / 2.04 TB/s & & data sheet\\
\bottomrule
\end{tabular}
\end{table}

\section{Reproducibility}\label{app:repro}

The complete results of this report are regenerated from the project root by the
following commands:
\begin{verbatim}
python3 sim/run_all.py            # 37 s -> results/*.json, *.csv, numbers.tex
Rscript figs/src/f01.R            # figure 01
Rscript figs/src/figs_data1.R     # data figures 12-17
Rscript figs/src/figs_data2.R     # data figures 08, 10, 21, 22, 25, 26
bash scripts/build-figures.sh     # schematic figures 02-07, 09, 11, 18-20, 23, 24
bash scripts/build-report.sh      # -> ACPU-1_Technical_Report.pdf
\end{verbatim}
The numerical experiments use fixed seeds, the PRNG is a PCG64 generator, and the model has no
randomness outside the Monte-Carlo study (which uses seed 42). The report is compiled with Tectonic (XeTeX) by \code{scripts/build-report.sh}. Figures are vector PDFs generated either by R (data) or by
Chromium (schematics); the UI figure is a set of screenshots of the running application.

\clearpage
\section{Acronyms}\label{app:acronyms}

\begin{multicols}{2}\small
\begin{itemize}[leftmargin=0pt,itemsep=0pt]
\item ACPU --- Atmospheric Calculation Processor Unit
\item AVT --- Atmospheric Vector Tile
\item AVISA --- ACPU Vector Instruction Set Architecture
\item BSPDN --- Backside Power Delivery Network
\item DA --- Data Assimilation
\item DVFS --- Dynamic Voltage and Frequency Scaling
\item EnKF --- Ensemble Kalman Filter
\item FIT --- Failures In Time (per $10^9$ device-hours)
\item FTLE --- Finite-Time Lyapunov Exponent
\item GAA --- Gate-All-Around
\item HBM --- High Bandwidth Memory
\item NoC --- Network-on-Chip
\item NWP --- Numerical Weather Prediction
\item SPPT --- Stochastically Perturbed Parameterisation Tendencies
\item SKEB --- Stochastic Kinetic Energy Backscatter
\item SRAM --- Static Random Access Memory
\end{itemize}
\end{multicols}
